\documentclass[11pt]{article}
\usepackage{fontspec}
\usepackage[french,english]{babel}
\usepackage{amsmath}
\usepackage{amssymb}
\usepackage{amsthm}
\usepackage{mathrsfs}
\usepackage{tikz-cd}
\usepackage{enumitem}
\usetikzlibrary{babel}

\newcommand{\printedpage}[1]{\marginpar{\footnotesize p.~#1}}
\newcommand{\msp}[1]{\marginpar{\footnotesize [#1]}}
\newcommand{\Ob}{\mathrm{Ob}\,}
\newcommand{\Hom}{\mathrm{Hom}}

\setlength{\parindent}{0pt}

\begin{document}

% ===================== page 11 =====================
\printedpage{53}
% head: A. Grothendieck : Dérivateurs, Ch. XII, Édition des Universités de Montpellier II et Paris VII

\msp{page 74}
But the inclusion of an initial object is always aspheric, and remains so under product with any object (i.e. $u'\backslash U$) of Cat (and even under any smooth base change \dots), whence the conclusion.

In what follows, we are mainly interested in the situation of amalgamated sums (where we change notation a little, now)
\[
\begin{tikzcd}
U \arrow[r, hook, "i"] \arrow[d, "p"'] & X \arrow[d, "q"'] & T \arrow[l, hook', dotted, "j"'] \arrow[d, dotted] \\
U' \arrow[r, hook, "i'"] & X' & T' \simeq T. \arrow[l, hook', dotted]
\end{tikzcd}
\tag*{$(*)$}
\]

The closed subcategory $T'$ complementary to $U'$ in $X'$ is then canonically isomorphic to the closed subcategory $T$ complementary to $U$ in $X$, the isomorphism being induced by $q$. Often we shall identify $T'$ and $T$, which is legitimate on viewing $X'$ as obtained from the gluing of $U'$ and $t$ \texttt{[rather }$T$\texttt{]}, by the bifunctor $U'^{\circ}\times T\longrightarrow \mathrm{Ens}$ made explicit in formula $(**)$ p.~72. We are mainly interested in isolating conditions which ensure the following two implications

\begin{itemize}
\item[a)] $i\in W\Longrightarrow i'\in W$ (and even $\Longleftrightarrow$).
\item[b)] $p\in W\Longrightarrow q\in W$ (and even $\Longleftrightarrow$).
\end{itemize}

\msp{page 75}
\textbf{Proposition 4.} \textit{Let $t\in\Ob T$. For $t\in\Ob\bar{U}$ to hold, it is necessary and sufficient that $q(t)\in\Ob\overline{U'}$. In other words, the ``boundary'' $\dot{T}=T\cap\bar{U}$ of $T$ in $X$ is ``the same'' as its boundary in $X'$.}

\textsc{Proof}. To say that $t\in\Ob\bar{U}$ also means
\[
U/t\neq\emptyset
\]
and implies, by continuity, that $q(t)\in\overline{q(U)}\subseteq\overline{U'}$. Conversely, if $q(t)\in\overline{U'}$, i.e. if there exists $u'\in U'$ such that
\[
\Hom(u',q(t))(=\pi_0(u'\backslash U/t))\neq\emptyset,
\]
it follows that $u'\backslash U/t\neq\emptyset$, a fortiori $U/t$ (which receives $u'\backslash U/t)\neq\emptyset$, q.e.d.

\textbf{Proposition 5.} \textit{Let}
\[
\begin{tikzcd}
Y \arrow[r, "i"] \arrow[d, "p"'] & X \arrow[d, "q"] \\
Y' \arrow[r, "i'"] & X'
\end{tikzcd}
\]
% foot: Version 7 décembre 2002
% ===================== page 12 =====================
\printedpage{54}
% head: A. Grothendieck : Dérivateurs, Ch. XII, Édition des Universités de Montpellier II et Paris VII
\textit{be a cocartesian square in \textup{Cat}. If $i$ is an open immersion (resp. closed), then $i'$ too, and the square is also cartesian. Ditto for $p$ and $q$. If $i$ and $p$ are open immersions (\footnote{\textbf{NB} We use \textit{``open immersion''}, \textit{``closed immersion''} in the \textit{strict sense}\,: it is a fully faithful \textit{monomorphism} whose image is a subcategory that is \textit{full,} either open or closed.}) (resp. \dots), then the situation is isomorphic to that of an object $X'$ of \textup{Cat}, of two open (resp. closed) subcategories $X$, $Y'$ of $X$ having as union $X'$, as intersection} % [sic] de $X$ -> de $X'$, X and Y' are opens of X' (their union is X')
\textit{$Y$. Conversely, in such a situationm the corresponding square is cocartesian, i.e. $X$} % [sic] situationm -> situation, (comma), typing error
\texttt{[rather }$X'$\texttt{]} \textit{is identified with the amalgamated sum}
\[
X'\simeq X\amalg_{Y}Y'.
\]

Thus we also have proposition 15, page 158, and its corollaries.
% foot: Version 7 décembre 2002

% ===================== page 13 =====================
\printedpage{55}
% head: A. Grothendieck : Dérivateurs, Ch. XII, Édition des Universités de Montpellier II et Paris VII

\section*{2\quad Amalgamated sums and $W$-equivalences.}

\subsection*{2.1\quad Cocartesian square axiom and cube axiom (or ``Five Lemma'').}

We now suppose given a basic localizer
\[
W\subseteq\mathrm{Fl}(\mathrm{Cat}),
\]
so $W$ satisfies W(1, 2 bis, 3). We do not assume for the moment for the other standard axioms W(4) (fibration axiom), W(5) (connectedness axiom), W(6) (direct limit axiom). But we are now going to introduce two additional stability axioms, of a different nature, concerning amalgamated sums.

\begin{description}[leftmargin=5em, style=nextline]
\item[W(7 bis)] (Cocartesian square axiom.) Let $X$ be an object of Cat, $U$, $U'$ two open subcategories of $X$ such that $U\cup U'=X$, let us consider the corresponding cocartesian diagram (cf. prop. 5 in the margin of the page \texttt{[above]}):

\msp{page 76}
\[
\begin{tikzcd}[column sep=small]
& U\cap U' \arrow[dl, hook', "i'_0"'] \arrow[dr, hook, "i_0"] & \\
U \arrow[dr, hook, "i"'] & & U' \arrow[dl, hook', "i'"] \\
& X &
\end{tikzcd}
\tag*{$(*)$}
\]

Then
\[
i_0\in W\qquad\Longleftrightarrow\qquad i\in W
\]
(and symmetrically
\[
i'_0\in W\qquad\Longleftrightarrow\qquad i'\in W).
\]
\end{description}

\textbf{NB} By duality, this axiom is equivalent to the similar axiom, with \textit{closed subcategories} $T$, $T'$ of $X$ with union $X$.

\textbf{Corollary.} \textit{Assume \textup{W(7 bis)}. Consider the following conditions on the preceding square.}

\begin{itemize}
\item[\textit{a)}] \textit{$U\cap U'$ $W$-aspheric.}
\item[\textit{b)}] \textit{$U$, $U'$ $W$-aspheric.}
\item[\textit{c)}] \textit{$X$ $W$-aspheric.}
\end{itemize}

\textit{Then (a and b) implies c, and (b and c) implies a.}

\textbf{NB} Morally, (c and a) ought also to imply b, i.e. if $U\cap U'$ and $X$ are $W$-aspheric, the same is true of $U$ and $U'$. But it does not seem that the cocartesian square axiom allows one to prove this. On the other hand, the axiom which follows will imply the preceding corollary in its complete form.
% foot: Version 7 décembre 2002

% ===================== page 14 =====================
\printedpage{56}
% head: A. Grothendieck : Dérivateurs, Ch. XII, Édition des Universités de Montpellier II et Paris VII

\begin{description}[leftmargin=5em, style=sameline]
\item[W(7)] (``Cube axiom'', or ``Five Lemma''.) Let there be a commutative cubical diagram
\[
\begin{tikzcd}[row sep=small, column sep=small]
U_1\cap U'_1 \arrow[rr] \arrow[dr] \arrow[dd, "f_{U_1\cap U'_1}"'] & & U_1 \arrow[dr] \arrow[dd, "f_{U_1}"{swap,pos=.25}] & \\
& U'_1 \arrow[rr, crossing over] & & X_1 \arrow[dd, "f"] \\
U_2\cap U'_2 \arrow[rr] \arrow[dr] & & U_2 \arrow[dr] & \\
& U'_2 \arrow[rr] \arrow[from=uu, crossing over, "f_{U'_1}"{pos=.25}] & & X_2
\end{tikzcd}
\]
representing a morphism (vertical arrows) from a

\msp{page 77}
cocartesian square\quad $\vcenter{\hbox{\begin{tikzcd}
U_1\cap U'_1 \arrow[r] \arrow[d] & U_1 \arrow[d] \\
U'_1 \arrow[r] & X_1
\end{tikzcd}}}$\quad of the type considered in W(7 bis) ($X_1$ union of two open subcategories $U_1$, $U'_1$) into another (with subscripts 2). Consider the following conditions.

\begin{itemize}
\item[a)] $f_{U_1\cap U'_1}\in W$.
\item[b)] $f_{U_1}$ and $f_{U'_1}$ are \texttt{[in]} $W$.
\item[c)] $f\in W$.
\end{itemize}

Then any two of these conditions imply the third. In other words, when one of these conditions is satisfied, the other two are equivalent to each other.
\end{description}

\textbf{NB} This axiom is equivalent to the dual axiom, in terms of objects $X_1$, $X_2$ each represented as a union of two closed subcategories.

\textbf{Corollary.} \textit{Assume \textup{W(7)} satisfied. Then the corollary to \textup{W(7 bis)} (an axiom which we do not have to use here), in its completed form (cf. the NB following it) is valid: any two of the conditions a), b), c) imply the third.}

It suffices to look at the unique morphism from the diagram $(*)$, p.~76, into the trivial diagram\quad $\vcenter{\hbox{\begin{tikzcd}[column sep=small]
& e \arrow[dl] \arrow[dr] & \\
e \arrow[dr] & & e, \arrow[dl] \\
& e &
\end{tikzcd}}}$\quad corresponding to the case where $X$ is the final category $e$, and $U$, $U'$ equal to $X$. Then the corollary is nothing other than W(7) applied to this particular case.

This argument makes one feel that W(7) is probably stronger than W(7 bis). Prompted by this feeling, I notice with pleasure that it is well founded:
% foot: Version 7 décembre 2002

% ===================== page 15 =====================
\printedpage{57}
% head: A. Grothendieck : Dérivateurs, Ch. XII, Édition des Universités de Montpellier II et Paris VII

\msp{page 78}
\textbf{Proposition 6.} \textit{The axiom \textup{W(7)} (five lemma) implies \textup{W(7 bis)} (cocartesian square axiom).}

We look at the following cubical diagram
\[
\begin{tikzcd}[row sep=small, column sep=small]
U\cap U' \arrow[rr, "i_0"] \arrow[dr, "i'_0"] \arrow[dd, "f_{U\cap U'}=i_0"'] & & U' \arrow[dr, "i'"] \arrow[dd, "f_{U'}=\mathrm{id}"{swap,pos=.25}] & \\
& U \arrow[rr, crossing over, "i"] & & X \arrow[dd, "f=\mathrm{id}_X"] \\
U' \arrow[rr, "\mathrm{id}"] \arrow[dr, "i'"] & & U' \arrow[dr, "i'"] & \\
& X \arrow[rr, "\mathrm{id}"] \arrow[from=uu, crossing over, "f_U=i"{pos=.25}] & & X,
\end{tikzcd}
\]
from the given cocartesian square into the cocartesian square obtained from $X$ and its two open subcategories $X$, $U'$, with intersection $U'$. Here the two vertical arrows $f$ and $f_{U'}$ are isomorphisms, hence \texttt{[in]} $W$. With the notation of W(7), we therefore have c) ($f\in W$), hence a) and b) are equivalent, i.e. $i_0\in W$ and $i\in W$ are equivalent.

\hfill q.e.d.

\textbf{Theorem 2.} \textit{The basic localizer $W_\infty$ satisfies the axiom \textup{W(7)}.}

This is surely a consequence of the cohomological criterion of \textsc{Artin-Mazur}, taking into account the Mayer-Vietoris exact sequence for both of the two cocartesian squares, and the noncommutative variants of Mayer-Vietoris. I have not written out the verification in due form, and should probably repair this omission.

\msp{page 79}
\textbf{Remark.} The dual version (in the sense of \textsc{Quillen}) of the cocartesian square axiom W(7 bis), i.e. the corresponding axiom for cartesian squares, is not valid for $W_\infty$. The relevant statement would be the following: let
\[
\begin{tikzcd}
X & X' \arrow[l, "q"'] \arrow[d, "f'"] \\
S & S' \arrow[l, "p"']
\arrow[from=1-1, to=2-1, "f"']
\end{tikzcd}
\]
be a cartesian square, with $f$ a \textit{$W$-fibration.} Then $f\in W\Longleftrightarrow f'\in W$. If $W$ satisfies W(4c) (which is the case for $W_\infty$), then $f\in W$ does indeed imply that its fibers are $W$-aspheric, hence that it is a universal $W$-equivalence, i.e. $f'\in W\ \forall\, p:S'\longrightarrow S$. But for the reverse implication $f'\in W\Longrightarrow f\in W$, one must evidently assume that
\[
\pi_0(p):\pi_0(S')\longrightarrow\pi_0(S)
\]
% foot: Version 7 décembre 2002

\end{document}
